% Reader-facing proof roadmap for the corrected Version 2 manuscript.
% This file is not included by main.tex automatically.  Insert only after the
% Section VIII closure theorem has been checked and the theorem statement has
% been updated consistently.

\section{Proof architecture}\label{sec:proof-architecture}

The proof has three logically separate components.  First, we compare the
ordinary-colouring and signed four-size first-moment roots.  Second, we prove a
normalised second-moment estimate for the signed profile.  Third, we amplify the
resulting positive-probability signed seed to a high-probability cocolouring.
We record the dependency structure before entering the technical estimates.

\subsection{The two first-moment locations}

Let $r_+(n)$ denote the continuous root of the ordinary profile first moment,
and let $r_4^{\mathrm{co}}(n)$ denote the root of the signed four-size
objective.  The four permitted deficits are $2,3,4,5$.  If $\delta_n$ is the
phase parameter and
\[
  A_4(\delta):=\log 2-D_4(\delta),
\]
then the root comparison gives
\[
  r_+(n)-r_4^{\mathrm{co}}(n)
  =
  \left[
    \frac{(\log2)^2}{4}A_4(\delta_n)+o(1)
  \right]\frac{n}{(\log n)^3}.
\]
We place the signed seed at
\[
  k_{\mathrm{co}}
  =
  \left\lceil
    \frac{r_4^{\mathrm{co}}(n)+r_+(n)}2
  \right\rceil.
\]
Thus the deterministic distance from the ordinary root is
\[
  r_+(n)-k_{\mathrm{co}}
  =
  \left[
    \frac{(\log2)^2}{8}A_4(\delta_n)+o(1)
  \right]\frac{n}{(\log n)^3}.
\]
The exact entropy certificate used below implies
\[
  A_4(\delta)>\log\!\left(\frac{1000}{639}\right)
\]
uniformly in the phase.

\subsection{The signed overlap identity}

For two ordered signed profile partitions, write $r=(r_{ab})$ for their overlap
table and let $H_r$ be the bipartite graph of cells with multiplicity at least
two.  Summing over the two sign vectors gives the exact identity
\[
  \operatorname{wt}_{\mathrm{sign}}(r)
  =
  \left(\prod_{a,b}g(r_{ab})\right)2^{\beta(H_r)},
\]
where $\beta(H_r)$ is the binary cycle rank.  We separate the normalised
second moment into partial diagonals, canonical high cells, and the residual
attachment.

The partial-diagonal contribution is handled directly.  The positive support
of the canonical high cells is a block matching, because every row and column
margin is at most the phase cap whereas each high cell exceeds one half of that
cap.

\subsection{High cells as a matching with deficits}

Fix a block matching $P$.  For $e\in P$, let $s_e,t_e$ be the endpoint block
sizes,
\[
  m_e:=\min\{s_e,t_e\},\qquad
  j_e:=m_e-h_e,
\]
and put
\[
  J:=\sum_{e\in P}j_e,\qquad H:=\sum_{e\in P}h_e.
\]
After summing the literal partial-stub-matching fibre, the exact aggregate
weight is
\begin{equation}\label{eq:v2-aggregate-high-weight}
  w(P,j)
  =
  \frac{\displaystyle
    \prod_{e\in P}(s_e)_{j_e}(t_e)_{j_e}}
  {\displaystyle
    (n)_J\prod_{e\in P}j_e!}
  \prod_{e\in P}g(j_e).
\end{equation}
No full completion is selected objectwise.  Equation
\eqref{eq:v2-aggregate-high-weight} is obtained by summing the entire finite
physical fibre.

For endpoint sizes $m,m+d$, define
\[
  R_{m,d}(h)
  :=
  \frac{\binom mh}{(d+1)(d+2)\cdots(d+h)}
  2^{-hm+h(h+1)/2}.
\]
The only global denominator change is
\[
  \frac{(n)_{J+H}}{(n)_J}
  =(n-J)_H\le n^H.
\]
Consequently
\begin{equation}\label{eq:v2-global-deficit-comparison}
  \frac{w(P,m-h)}{w_{\mathrm{full}}(P)}
  \le
  \prod_{e\in P}n^{h_e}R_{m_e,d_e}(h_e).
\end{equation}

The canonical high condition implies $2h_e<m_e$.  Hence
\[
  n^hR_{m,d}(h)
  \le
  \left(\frac{nm}{2^{\lfloor2m/3\rfloor}}\right)^h.
\]
For the four phase sizes,
\[
  \rho_n
  :=
  \max_m\frac{nm}{2^{\lfloor2m/3\rfloor}}
  =O\!\left(\frac{(\log n)^{7/3}}{n^{1/3}}\right)=o(1).
\]
Thus the sum of all positive deficits in one selected cell is at most
$2\rho_n$ for large $n$, and the entire deficit fibre contributes
\[
  \exp\!\left\{
    O\bigl(n^{2/3}(\log n)^{4/3}\bigr)
  \right\}.
\]
The endpoint transportation estimate contributes
$\exp\{O(\sqrt{n\log n})\}$.  Together with the partial-diagonal estimate this
gives
\[
  \operatorname{BareSkeletonSum}_n
  \le
  \exp\!\left\{
    o\!\left(\frac{n}{(\log n)^4}\right)
  \right\}.
\]

\subsection{Residual attachment by injective restriction}

Let $M$ be the exposed high-support matching.  Deletion of $M$ is injective on
the family of even residual edge sets: the symmetric difference of two
completions would be an even subset of a matching and is therefore empty.  For
nonnegative activities $(q_e)$,
\[
  \sum_{F\ \mathrm{even}}
    \prod_{e\in F\setminus M}q_e
  \le
  \prod_{e\notin M}(1+q_e).
\]
The local increment activity is bounded pointwise by the same $q$ activity.
The intrinsic and complementary residual regimes therefore yield
\[
  \operatorname{AttachmentSum}_n
  \le
  \operatorname{BareSkeletonSum}_n
  \exp\!\left\{
    \eta_n\frac{n}{(\log n)^4}
  \right\},
  \qquad \eta_n\longrightarrow0.
\]
This proves the required normalised second-moment estimate.

\subsection{Amplification and the final intersection}

Paley--Zygmund supplies a positive-probability signed seed.  A simultaneous
leftover-colouring lemma and a one-Lipschitz cocolourable-capacity variable,
combined with bounded differences, upgrade this seed to a high-probability
upper bound for $\zeta(G_n)$ with additive loss
$o(n/(\log n)^3)$.  The ordinary first-moment argument gives the corresponding
high-probability lower bound for $\chi(G_n)$.  Intersecting the two events by a
union bound completes the proof.
